
On
suppose toujours que \(a\in K\). La multiplication pour la loi \(\circ\) par \(\lambda_t(a)\) définit un morphisme de foncteurs : \(\Ghat_K\to\Ghat_K\), caractérisé par sa valeur pour l'élément \(1+Xt\) de \(\Ghat(K[X])\). On a donc :
\[
  \lambda_t(a)\circ(1+Xt)=\lambda_{Xt}(a)\qquad\text{d'après (2.3.2), donc :}
\]
\begin{equation}\tag{2.9.3}
  \lambda_t(a)\circ(1+Xt)
  =(1+Xt)^a
  =\exp(a\Log(1+Xt))
  =\sum_{n=0}^{\infty}\binom{a}{n}X^nt^n .
\end{equation}
On en déduit :
\begin{equation}\tag{2.9.4}
  \lambda_t(a)\circ\varphi=\varphi^a
  \quad(=\exp(a\Log\varphi)) ;
\end{equation}
En effet, les deux membres sont des foncteurs additifs égaux pour \(1+Xt\). De (2.9.3) et 1.15, il résulte que le \(n\)-ième polynôme universel définissant la multiplication par \(\lambda_t(a)\) est isobare de poids \(n\)~; il sera noté \(R_{n,a}\).

Notons enfin que si \(A\) est une \(K\)-\(\lambda\)-algèbre, i.e. un \(\lambda\)-anneau muni d'un \(\lambda\)-homomorphisme : \(K\to A\), et si \(b\in A\), on a :
\begin{equation}\tag{2.9.5}
  \lambda_t(ab)=(\lambda_t(b))^a .
\end{equation}
Ceci résulte en effet de \(\lambda_t(ab)=\lambda_t(a)\circ\lambda_t(b)\), et de (2.9.4).

\section*{3. Les opérations \(\gamma\)}

On se place encore sur la catégorie \(\mathcal C_0\).

3.1. Soit \(K\) un \(\lambda\)-anneau. Nous allons définir à partir des opérations \(\lambda^i\) de \(K\) de nouvelles oérations, notées \(\gamma^i\). Elles sont destinées à introduire une notion de classes de Chern sous des conditions qui seront indiquées plus tard. Pour le moment, donnons une «~justification~» des définitions qui suivent : si \(\underline{E}\) est un fibré vectoriel de rang \(d\), la \(d\)-ième classe de Chern de \(\underline{E}\) devra être (quelle que soit la théorie des classes de Chern envisagée) la «~somme alternée des puissances extérieures 
